\chapter[Why Literature Retells Old Stories]{Why Literature Keeps Rewriting Old Stories}
\section{A Face Revised Countless Times}
First pick something up: a cinema ticket.

Suppose it is a ticket stub from some summer, for the film Ne Zha. The Nezha on screen probably bears no resemblance to the divine child you imagined: dark circles surround his eyes, sharklike teeth fill his grin, his hands stay in his pockets, and he walks with a careless swagger, speaking like a young thug who respects nobody. His declaration, "My fate is mine, not heaven's", was later printed on countless stationery products and posters.

Now ask people of different ages what Nezha looks like. Their answers will fight one another.

Your grandparents' Nezha probably comes from the 1979 animated film Nezha Conquers the Dragon King: a handsome boy in white who, rather than implicate his parents and ordinary people, cuts his own throat with a sword in pouring rain, saying as he dies that he returns his flesh and bones to his father. This tragic Nezha brings tears.

Turn further back, to the Ming novel Investiture of the Gods, and he changes again: he disturbs the sea and kills the Dragon King's third son, returns his bones to his father and flesh to his mother, then receives a new body of lotus flowers and roots from his master. After resurrection, astonishingly, he pursues his own father Li Jing with murderous intent, forcing Li Jing always to carry a protective pagoda. This Nezha rebels almost savagely.

Dig back as far as possible and you encounter something surprising to many Chinese people: Nezha was not originally a Chinese child at all. His name transliterates an ancient Indian name. Originally a Buddhist guardian deity, son of the northern Heavenly King Vaishravana, he travelled from India to China with Buddhism. Chinese novelists gradually stripped away his armour, gave him flesh and blood, and made him the third son of Chentang Pass's military commander.

One name, at least four completely different faces: tragic, savage, roguish, sacred, all Nezha. No "original author of Nezha" can stand up and declare everyone else's version wrong, because each version receives the story from its predecessor and passes it onward. That cinema ticket bears the latest version of a character rewritten for over a millennium.

Here is this chapter's question: why does literature keep rewriting old stories? Is rewriting laziness or creation? When an old story changes beyond recognition, is it still the same story?

\section{Tears in Bianjing's Entertainment Quarter}
Now come into a scene with me.

Time: the Northern Song, around the eleventh century CE. Place: a famous wazi entertainment quarter in the capital, Dongjing Bianliang, today's Kaifeng in Henan.

Imagine a huge open-air combination of amusement park and shopping street, open day and night, crowded with snack sellers, acrobats, singers, and wrestlers. Among its biggest attractions is the storyteller's pitch. "Speaking" here means not chatting but a specialised craft: historical tales, Buddhist narratives, fiction. One person, one mouth, and one fan can root hundreds of listeners to the spot all afternoon.

The performance you squeeze into today is "telling the Three Kingdoms". Silk-clad merchants mingle with resting porters and children their exasperated families have sent out. The storyteller strikes his wooden block; silence falls. He describes Liu Bei's defeat and desperate flight. Several children beside you frown; one's eyes redden. Later, Cao Cao suffers defeat and flees. The same red-eyed children suddenly clap, some jumping for joy.

This scene is not invented. The Northern Song writer Su Shi recorded an almost identical incident in Dongpo's Forest of Notes: families unable to control mischievous street children gave them coins to hear stories. "When the Three Kingdoms were told, hearing Liu Xuande defeated, they frowned, some shedding tears; hearing Cao Cao defeated, they immediately cheered with delight." Liu Bei's defeat made them frown and cry; Cao Cao's brought instant celebration.

Notice an easily missed fact. When Su Shi recorded this, the novel Romance of the Three Kingdoms did not exist. Luo Guanzhong would write it centuries later, around the Yuan--Ming transition. The Three Kingdoms that made these children cry and laugh therefore came not from a book they had read, but from generations of storytellers refining it performance by performance. Chen Shou's historical Records of the Three Kingdoms supplied dry accounts of people and events. Countless nameless storytellers turned its Liu Bei into embodied benevolence, Cao Cao into a white-faced schemer, and battles into gripping dramas of tears and laughter.

When Luo Guanzhong finally sat down to write Romance of the Three Kingdoms, he faced no blank page, but a mass of old stories told for centuries in entertainment quarters. He organised, selected, added, deleted, and told them again: plainly, he rewrote them. After him, Mao Zonggang and his father rewrote them, theatre companies rewrote them, and television, games, and comics are rewriting them today.

Remember that entertainment quarter. The rest of this chapter answers the questions hidden inside it.

\section{Who Owns a Story?}
Set out the conflict before rushing to answers.

That scene can be told in two entirely different ways.

Version one: a beautiful tradition. The Three Kingdoms belong to everyone. History supplied a skeleton; thousands of storytellers, theatre companies, novelists, and painters spent over a millennium giving it today's flesh and blood. Luo Guanzhong stole nobody's property because there was no single somebody: the story's author was "everyone".

Version two: centuries of rewriting revelry, with victims among the revellers. The historical Cao Cao was a statesman, commander, and poet capable of writing "With wine before us, let us sing: how long is human life?" Yet successive storytellers blackened his name until the stage reduced him to a treacherous minister resembling a white-nosed clown. Once that image solidified, nobody wanted to reopen his case. If Cao Cao knew, he would probably send every storyteller a solicitor's letter.

Behind these versions lie three real questions:

\textbf{Question 1: Is rewriting inheritance or parasitism?} When writers decline to invent stories from nothing and borrow old plots and characters instead, are they enriching tradition or hitching a free ride, sparing themselves the labour of invention?

\textbf{Question 2: Who has the right to rewrite?} Buddhist monks, Ming novelists, and modern animation directors have all rewritten Nezha. If someone tomorrow makes him utterly unacceptable to you, do you have grounds to say the adaptation is wrong? What grounds?

\textbf{Question 3: Is fidelity to the original a virtue or a shackle?} Every screen adaptation of a classic provokes arguments about "ruining the original" or "perfect fidelity". But wait: the novel Journey to the West, the "original" faithfully followed by its television adaptation, itself radically rewrites Xuanzang's real pilgrimage. What kind of fidelity is being faithful to a rewriter's rewriting?

None of these has a ready answer. They concern how literature works, how traditions form, and whether "originality" carries the weight we suppose. Before continuing, take a position: if a director transforms his childhood idol Monkey King beyond recognition, but makes a wonderful film, does he owe the older Journey to the West an apology?

\section[Originality and Tradition Make Their Cases]{Originality and Tradition Each Make Their Case}
Two voices have long argued over rewriting. This time, give each the fullest possible hearing before judging.

\textbf{Position 1: Originality comes first; rewriting needs boundaries.}

Present this case thoroughly. It deserves care because today's internet often mocks protecting original work as being "old-fashioned" or "copyright police", unfairly dismissing its concerns.

Its reasoning has at least three layers. First, \textbf{creation must earn a living}. Inventing a good story from nothing is expensive, risky labour; rewriting is cheaper. If rewriters can freely take another's central story, characters, and fictional world, change names, and profit, who will undertake the hardest, riskiest original work? Protecting originality protects the earnings of the first storyteller, encouraging more stories to exist. Second, \textbf{rewriting can hurt}. Characters are not building blocks: they connect to authors' labour and readers' feelings. Making a beloved childhood companion sleazy or vulgar causes real distress. Cross-cultural rewriting raises further risks: a powerful creator casually reshaping another people's story and widely distributing a distorted version does more than raise literary questions; it disrespects that people. Third, \textbf{excessive rewriting crowds out originality}. If markets offer only the hundredth remake of familiar stories, new ones may never even be noticed. The worry is not one adaptation, but a literary world containing nothing except adaptations.

This case is far from weak. Law supports it too: almost every country has copyright legislation reserving territory around original creators that others cannot casually enter.

\textbf{Position 2: Stories belong to everyone; literary history is a history of rewriting.}

This side holds an intimidatingly thick stack of evidence from across the world.

Start in Greece. The Iliad and Odyssey, revered in the West as literature's beginnings, bear Homer's name. Yet the "Homeric question", debated for over two millennia, asks precisely whether such a person existed. The mainstream view holds that folk singers performed the epics across many generations, improvising modifications and rearrangements each time, before they were fixed in writing. Western literature's first foundation stone is itself a building assembled through innumerable rewritings, without one sole author.

Next, Britain. Shakespeare is almost a synonym for originality. Yet inspect his plays: Hamlet adapts an old Danish legend; Romeo and Juliet an established Italian love story; King Lear and Macbeth have precursors in histories and earlier drama. His independently invented plots can almost be counted on one hand. His genius lay in telling old stories with unprecedented depth, not inventing them. Under the strictest originality-first test, literary history's greatest dramatist was a professional rewriter.

Then India and Southeast Asia. The Indian epic Ramayana tells how Prince Rama rescues his wife Sita. In Thailand it became the Ramakien: characters acquired Thai names and events absorbed Thai settings and customs. Today its scenes cover Bangkok's Grand Palace mural galleries; it is a national treasure, and Thai kings even use Rama as a royal title. Cambodia, Laos, and Indonesia each grew another version. Nobody accused Thailand of stealing India's story. Stories travelling with merchants and monks and taking root wherever they arrived were entirely ordinary then.

China offers equally strong evidence. Journey to the West did not spring fully formed from the head of Wu Cheng'en, its conventionally accepted author. Xuanzang's Tang journey to India was real history. Centuries later, The Story of How Tripitaka of the Great Tang Procured the Scriptures already supplied a "Monkey Pilgrim" as escort. Yuan dramas gradually brought Pigsy and Sandy aboard; folk storytelling embellished events for centuries before the Ming hundred-chapter novel appeared. Today's Journey to the West is the joint work of dozens of centuries, hundreds of nameless authors, and one great synthesiser. The supposed original is itself rewriting's destination---and the next rewriting's departure point.

This evidence does not make defenders of originality retreat, because the sides discuss somewhat different things: one concerns \textbf{authors earning a living now}, the other \textbf{literature surviving across long historical periods} through rewriting. The difficult question is where to draw a line. Before drawing one, we need a finer tool.

\section[Thinking Bridge: A Rewriting Spectrum]{A Thinking Bridge: Draw a Spectrum of Rewriting}
Arguments about reckless adaptation or plagiarism usually divide people between two poles: rewriting or copying. A finer ruler immediately clarifies many disputes. First a term: literary scholars call texts' continual dialogue with other texts \textbf{intertextuality}. No piece of writing is an island. Reading about Kong Rong yielding the larger pears to his brothers and recalling your own family's struggle over chicken drumsticks creates an echo between texts: intertextuality at work. Once we acknowledge it, the question ceases to be whether borrowing occurs---it always does---and becomes how. Our ruler is therefore a \textbf{spectrum}, arranging relationships with earlier texts from near to far:

\begin{quote}
Allusion $\rightarrow$ Quotation $\rightarrow$ Reworking $\rightarrow$ Parody $\rightarrow$ Adaptation $\rightarrow$ Continuation / fan fiction $\rightarrow$ Plagiarism
\end{quote}
Examine them individually. \textbf{Allusion} compresses an old story into a password. Xin Qiji's line "Lian Po has grown old: can he still eat heartily?" brings an elderly general's entire wish to fight into a poem in only seven Chinese characters. Those who recognise it immediately understand that the poet means himself. \textbf{Quotation} declares its source and imports another's sentence for a new purpose. \textbf{Reworking} dissolves old sentences into new ones like salt in water; recognition depends on readers' experience. \textbf{Parody} deliberately imitates a work's manner to make fun of it. Cervantes's Don Quixote imitates chivalric romances throughout: an impoverished country gentleman reads himself mad, mounts a skinny horse to perform heroic deeds, and attacks windmills mistaken for giants. After laughing, you discover that it effectively sends the whole genre to its grave while creating something greater than any chivalric romance. \textbf{Adaptation} moves an entire story into a new medium or age: novel into film, ancient tale into modern city. \textbf{Continuations and fan fiction} arise when readers refuse to let a story end and write onward themselves. \textbf{Plagiarism} occupies the far end: taking another's work while pretending it is yours.

The spectrum is ready; using it matters. Three questions locate a rewriting and reveal its nature:

\textbf{First: Does it acknowledge its source?} Allusion, quotation, and parody hope for, even depend on, recognition. Without it, the reference fails and parody loses its joke. Plagiarism fears precisely that recognition. The first distinction therefore concerns not the scale of alteration, but whether the author wants you to recognise or hopes you never discover. An unchanged sentence with attribution is open quotation; a radically transformed work claiming complete originality may still plagiarise.

\textbf{Second: What changes, and why?} To honour or contradict? To entertain or profit without effort? Du Mu, Wang Anshi, and Li Qingzhao each approach Xiang Yu with different concerns, as the next section shows. They alter the story's meaning rather than its events. A hastily produced cash-grab remake often changes only the actors' faces.

\textbf{Third: What does it let readers see anew?} Good rewriting shines a torch back onto the old text, revealing overlooked features. After Don Quixote, you cannot read chivalric romances as before: laughter mixes with sadness. Bad rewriting is merely a mirror reflecting its own laziness.

Next time arguments about an adaptation ruining a book or a song plagiarising another fill your feed, do not immediately choose sides. Ask: does it acknowledge its source, what is its purpose in changing it, and what new thing can I see? Most disputants turn out to be arguing about different questions.

\section[At the Wu River: Four Rewritings]{At the Wu River: Four Rewritings of One Scene}
Now read a little primary material. We will watch one historical moment acquire four completely different meanings across centuries.

The moment is Xiang Yu's death. In 202 BCE, near the end of the Chu--Han struggle, Liu Bang's armies surrounded him at Gaixia. Breaking out, he fled to the Wu River, where a local official brought a boat and urged him to cross home to Jiangdong and begin again. Sima Qian records his final answer in "Basic Annals of Xiang Yu", Records of the Grand Historian:

\begin{quote}
King Xiang laughed: "Heaven destroys me; why should I cross? I led eight thousand sons of Jiangdong west across the river, and now not one returns. Even if their fathers and brothers pity me and make me king, how could I face them?"
\end{quote}
In plain terms: if heaven wants me dead, what would crossing achieve? I brought eight thousand Jiangdong men west, and none returns alive. Even if their families pity me and still accept me as king, how could I face them? After speaking, Xiang Yu dismounted, fought on foot, killed hundreds, and cut his own throat.

Notice that this text is already the first rewriting. Decades separated Sima Qian from Xiang Yu; he could not have heard the riverside conversation, reconstructing it from reports and historical materials. More interesting is its placement: the Records reserves "Basic Annals" for emperors. Xiang Yu never became emperor and lost to Liu Bang, yet Sima Qian placed him alongside emperors in his own Basic Annals. The historian's pen rewrites and reassesses, rejecting the formula that victors become kings and losers bandits.

Then the poets take their turns.

Over eight centuries later, the Tang poet Du Mu passed the Wu River Pavilion and wrote "Inscribed at the Wu River Pavilion":

\begin{quote}
Victory and defeat are unpredictable in war; a true man bears shame and humiliation. Jiangdong has many talented sons; who knows whether he might have returned to fight again?
\end{quote}
The sense: victory and defeat are ordinary military affairs; endurance of shame proves manhood. Jiangdong has plenty of talent; regrouping and returning might have produced another outcome. Du Mu reads wasted possibility: Xiang Yu could not bear losing and rashly threw away his chance to recover. Himself ambitious yet frustrated, Du Mu uses Xiang Yu to say that a real hero must endure defeat.

A century or so later, Northern Song reformer Wang Anshi read Du Mu's poem and wrote "Another Inscription at the Wu River Pavilion" specifically to argue:

\begin{quote}
A hundred battles leave brave men weary and grieving; one defeat in the Central Plain is hard to reverse. Even if Jiangdong's sons remain today, would they return to fight beside their king?
\end{quote}
The sense: years of war have exhausted and embittered the soldiers; after defeat in the Central Plain, the larger situation is irretrievable. Even surviving Jiangdong men might refuse to follow him again. Wang directly contradicts Du: it is not that Xiang Yu refuses endurance, but that endurance cannot help. Years of fighting have scattered allegiance; timing, terrain, and popular support all oppose him. Jiangdong will offer no second chance. A political reformer, Wang habitually studies the overall situation and structural conditions, not simply an individual's capacity to endure.

Another century and more brings Southern Song poet Li Qingzhao. She experienced the Jingkang catastrophe: Jin armies invaded, the Northern Song fell, and its court abandoned northern territory to flee across the Yangtze. She wrote "Summer Quatrain":

\begin{quote}
In life, be a hero among people; in death, a hero among ghosts. Still I think of Xiang Yu, who would not cross to Jiangdong.
\end{quote}
The sense: be outstanding in life and heroic even in death. I remember Xiang Yu because he refused to cross and cling ignobly to life. Ostensibly praising him, every word stabs her own age: the court and all its officials crossed south, and felt entirely comfortable doing so. His refusal condemns their willingness.

One scene, four readings. Sima Qian finds a tragic hero's dignity; Du Mu the possibility of enduring humiliation; Wang Anshi the cruelty of irreversible circumstances; Li Qingzhao anger at dishonourable survival. The event never changes: Xiang Yu did not cross. What that refusal means receives an answer suited to each generation's time.

Here is rewriting's central act: \textbf{the old story is an empty stage onto which each generation brings its own questions.} Du asks how to face defeat; Wang how to understand circumstances; Li how to face national humiliation. They do care about Xiang Yu, but they also need him: a hero everyone recognises, capable of giving their own unspoken thoughts weight and clarity.

\section[Why Old Stories Persist: Three Accounts]{Why Literature Needs Old Stories: Three Explanations}
We can now offer genuinely different answers to our central question. First separate four things: what happened, literature across peoples repeatedly rewriting old stories for millennia, belongs to evidence and is scarcely disputed; why it happened belongs to interpretation, where argument begins; how participants understood it, Li Qingzhao regarding herself as criticising the present through the past rather than "rewriting", belongs to their own account; whether rewriting is good belongs to value judgement. We compare the second layer.

\textbf{Explanation 1: Old stories are preinstalled codes: cognition and efficiency.}

A major storytelling cost is orienting listeners: who someone is, why they act, what rules govern their world. Old stories preinstall a code in readers' minds. Mention Monkey King and a Chinese reader instantly sees the golden-hooped staff, tightening-headband spell, seventy-two transformations, and unruly temperament. No explanation needed; the author starts work immediately. Better still, readers' expectations let \textbf{rewriting generate power through a gap}: expecting Nezha as an obedient child and seeing a swaggering demon child creates surprise and comedy. Without old stories underneath, the new story cannot even be unexpected, because nothing was expected. This explanation neatly accounts for the explosive potential of adapting classics. Its limits: it cannot explain choosing obscure texts, or the gulf between rewritings using identical codes, one a masterpiece, another merely chasing attention.

\textbf{Explanation 2: Old stories are low-risk products: commerce.}

Now borrow investors' eyes rather than readers'. A film can cost hundreds of millions: gamble on an unknown story or one everyone recognises? The answer seems obvious. Hence Hollywood's yearly remakes of fairy tales, comics, and older films, and companies' frantic pursuit of adaptation rights to online novels and classics. An existing audience is like having half the tickets sold in advance. Sharp and candid, this explains our age's abundance of remakes, sequels, and reboots. Its limit: if everything concerns money, fan fiction becomes inexplicable. Millions worldwide earn nothing while spending nights writing sequels and drawing derivative art for fellow fans, donating time and money. Commercial accounting cannot balance that ledger. Other motives must exist.

\textbf{Explanation 3: Rewriting contests the right to tell: politics.}

Stand further back and ask who wrote old stories: mostly winners, powerful people, literate people. The defeated, weak, humiliated, and injured often appear as villains, jokes, or scenery. Rewriting lets later people, especially descendants of those once denied a voice, reclaim narration. A villain's prequel reveals an inner journey; a woman speaks words the original male author withheld; colonised peoples retell the coloniser's textbook heroic epic to reveal its bloody underside. Cao Cao's progressive blackening in that entertainment quarter offers the reverse lesson: controlling narration can flatten a complex person into a white-painted face. This account illuminates what the others miss, including the fierce anger of some adaptations. Its limit: turning all rewriting into power struggles overlooks another equally real motive. Many rewrite simply because they love a story until their fingers itch. Calling all affection calculation is itself a misreading.

Each explanation holds some truth. Codes explain rewriting's usefulness; commerce its abundance; power its sharpness. A single adaptation often contains all three forces.

\section[Further Challenge: One Hero, Many Faces?]{A Deeper Challenge: Are All Heroes the Same Person?}
Now introduce our weightiest theory, taking rewriting to its extreme: perhaps only one story exists.

American scholar Joseph Campbell proposed this. His 1949 The Hero with a Thousand Faces compared myths, epics, and religious narratives worldwide: Greek Odysseus, India's Buddha, the biblical Moses, Icelandic heroic sagas. Beneath wildly different appearances, he announced, their skeletons are astonishingly similar. He called the shared structure the "hero's journey" or "monomyth". Its rough route is:

\begin{quote}
Ordinary world $\rightarrow$ Call to adventure $\rightarrow$ Refusal $\rightarrow$ Meeting the mentor $\rightarrow$ Crossing the threshold $\rightarrow$ Trials and ordeals $\rightarrow$ Innermost cave $\rightarrow$ Desperate struggle $\rightarrow$ Winning the treasure $\rightarrow$ Return $\rightarrow$ Bringing something beneficial back to the ordinary world
\end{quote}
Consider Journey to the West: Tripitaka receives the mission to fetch scriptures, the call; Monkey King awaits his master beneath Five Elements Mountain, mentor and threshold; eighty-one ordeals supply trials; true scriptures finally return to benefit the eastern land, the treasure brought home. Or Harry Potter: the ordinary boy beneath the stairs receives an owl's letter, the call; Hagrid leads him into Diagon Alley, crossing the threshold; Hogwarts supplies trials; Voldemort finally confronts him, the innermost cave. Star Wars works too, unsurprisingly: George Lucas publicly acknowledged closely studying Campbell while developing it, building his film around this skeleton.

The theory immediately reveals two sides.

\textbf{Its usefulness is real.} It gives viewers a skeletal map, letting them anticipate many films: departure, setbacks, rock bottom, recovery. You roughly know what comes next. Writers receive a construction plan too. Hollywood turned it into a production manual, with screenwriting guides specifying the page for the call and for hitting bottom before rebounding. Commercial films worldwide consequently resemble one another because they really share a blueprint.

\textbf{Its dangers are equally real, with at least three layers.} First, description becomes recipe. Campbell originally observed similarities; once the plan exists, production becomes replication. When every film rebounds at the same moment, audiences yawn early: discovery deteriorates into formula. Second, fitting one mould flattens genuine differences. Many myth scholars criticise Campbell for selecting and editing favourable passages from actually diverse traditions to force one framework. Your universal law may be a target painted around an arrow already fired. Third, more subtly, the mould becomes a judge. Treat the hero's journey as the standard for good stories and other forms become failures at storytelling. Yet Chinese chapter novels value introduction, development, turn, and conclusion; Japanese Noh scarcely advances a plot; many peoples' epics lack a solitary hero overcoming successive obstacles. Such works have been great within their traditions for centuries.

A warning sign belongs here. First hearing this theory, many exclaim: "Humanity's shared psychological code---a scientific discovery!" That is a misjudgement. The hero's journey is no law scanned out of human brains. A twentieth-century scholar in his study generalised from some myths he had read. Generalisation can illuminate, but is always limited by what he read, missed, and edited. Used as a searchlight, it illuminates a large area; used as a cage, it confines you.

Its final lesson may be its own unfinished sentence: all heroes perhaps travel one road, yet stories move us through each hero's particular manner of walking, never the road alone. Old skeleton, new flesh: that is perhaps the best definition of rewriting.

\section[Old Stories in the Remix Age]{When Rewriting Becomes Everyday: Old Stories in the Remix Age}
This ancient question now plays ten times louder inside your phone.

First, remixing. Every minute short-video creators rewrite: editing the Three Kingdoms television drama into workplace tutorials, redubbing the four great Chinese novels' characters with modern lines, grafting new endings onto classic films. Guichu remix videos break and recombine old footage to make it say what the original never said. Parody through editing software strikingly resembles Cervantes's pen attacking chivalric romances four centuries ago. Scale differs: there was one Don Quixote; hundreds of millions now wield editing software.

Next, fan works. Countless enthusiasts write sequels, draw derivative pictures, and film little scenes for beloved stories, without payment, simply from love. Commerce alone cannot explain this immense creative flood. It also repeatedly strikes walls: rights holders enforce claims, platforms remove work, fans bitterly dispute what respect for an original requires. Our spectrum becomes a weapon and shield every day here.

Then the century-long adaptation argument. Each classic remake automatically splits comments between "ruining the original" and "freedom to adapt". Here lies a recurring misjudgement: \textbf{fidelity to the original has never been the standard for good adaptation}. Word-for-word, step-for-step screen versions often seem as lifeless as textbook recitation, because novels and films speak different languages; translation necessarily gains and loses. Ne Zha, accused of radical distortion, transformed a tragic boy into a roguish demon child but gave hundreds of millions a resonant new meaning. Judge not resemblance, but our three questions: does it acknowledge its source, why does it change it, and what more does it let us see?

A sharper new question is cultural appropriation. When American animation companies turn China's Mulan, Pacific islanders' legends, and Africa's Lion King stories into profitable global blockbusters, their stories' homelands receive little. Sometimes distortion is sold to global audiences as authenticity. Across cultural boundaries with unequal power, "stories belong to everyone" changes flavour: the speaker is often the one strong enough to take them. Distinguish several things. Such disputes repeatedly occurring worldwide is established fact. Unequal power and rewards, rather than adaptation itself, plausibly explain the problem. Whether a particular film counts as appropriation remains an open question requiring individual argument.

The newest variable is AI. Chatbots have read immense quantities of humanity's millennia of writing and can "rewrite" a passage in seconds. They have no love, anger, or impulse to vindicate someone: none of this chapter's three motives, yet they rewrite faster than anyone. Readers or authors? What do they owe the original authors whose collective stories taught them? Lawyers argue and writers protest; nobody yet knows the answers. Your generation will necessarily help supply them.

\section[Your Turn: What Does a Rewriter Owe?]{Your Turn: What Does a Rewriter Owe the Original?}
Return to the opening cinema ticket.

The dark-eyed Nezha shares a name with the Buddhist guardian, the father-pursuing youth in Investiture of the Gods, and the white-clad boy dying in 1979's rain. Every rewriting connecting them receives applause and causes somebody pain.

Now your turn. Imagine a director thirty years hence taking up Nezha and making him unrecognisable: perhaps he rebels against nobody, becomes the villain, or moves to Mars. You sit before some screen unimaginable today, watching this grandchild of "your Nezha".

By what standard will you decide whether this legitimately carries on a tradition or unforgivably ruins it?

Before answering, lay out your tools: the originality camp's warning that unprotected sources eventually stop producing new stories; the traditionalists' evidence that Homer, Shakespeare, and Journey to the West all arise from rewriting; the spectrum and its three questions about source, intention, and new contributions; the four poems at the Wu River, one scene nurturing four sincere meanings; and Campbell's lesson that even the seductive theory of one universal story can become a cage.

Does a rewriter owe the original a debt? If so, does attribution repay it, or must the new work add something substantial? If an adaptation makes billions love an old story again but leaves the old version unread forever, is that rescue or murder?

There is no standard answer. But notice something small: when you begin answering, you already draw on this chapter's Nezha, Xiang Yu, Cao Cao, and Don Quixote. You use old stories to convey your own new meaning. From your first words, you too join this river flowing for millennia.
